\subsection{Validation and the decision to keep this optional}

The 32 focused tests pass in 3.991 seconds. They include forty small
arithmetic instances checked on exhaustive potential grids, malformed
optimality witnesses, disconnected incidence components, 20,000-bit scaled
heights, cancellation, strict option types and exact positive-proof budget
boundaries. The final maintained suite passes all 1,197 tests in
232.354 seconds. An earlier full run before changing the default allowance
also passed; the final run covers the disabled-by-default policy.

The source audit independently inspects every distinct root-extremal
surface on the 74 inputs with at most twelve crossings, together with each
original optimum. All 787 surfaces agree with Regina on connectedness,
orientability, Euler characteristic and compressing-disc status. There are
713 alternatives, and 21 of the 74 sources exhibit different Euler
characteristics among their tied optima. Full witnesses for the two named
counterexamples and the new positive are retained. On nineteen ordinary
portfolio inputs, both zero-root and four-root settings return the expected
verdicts without reaching the native stage.

The first audit reached the final ordinary-corpus loop but failed because
its driver used an incorrect corpus-helper attribute. The helper reference
was corrected and the failure log retained. The final pinned audit takes
103.135 seconds, overlapping correctness tests; this is not a performance
comparison. Its 485 runtime, driver, fixture and prior-audit hashes agree
with the final benchmark. The independent span and source-verifier modules
are byte-identical to the previous baseline \code{a8658ec40}.

Two serial five-round experiments compare the complete recognizer with
face search explicitly off or set to four roots. Each case has one warm-up
round and then five measured rounds, each randomly ordering off, off A/A,
on and on A/A calls. No tests, audits, builds or commits overlap the timed
runs. Earlier shortcuts are disabled so that the native stage is reached;
both arms have no local seed cap and retain the twelve-second global
recognition deadline. All requested tree trials remain the same in both
arms, including 24 for the late-tree control. A returned native proof is
additionally replayed and hashed inside the timed endpoint.

The first experiment motivated the zero default. The second repeats the
comparison with the final source pins. All 280 measured calls and 56
warm-ups across the two experiments finish; per-implementation verdicts
and proof hashes are stable across both runs. The final run alone has
140 measured calls and 28 warm-ups and takes 15.003 seconds. Its primary-arm
median times and median paired ratios are shown below; a median paired
ratio need not equal the ratio of median times.

\begin{center}
\small\begin{tabular}{lrrrrr}
\toprule
Input & Off (ms) & On (ms) & Off/on & Off A/A & On A/A \\
\midrule
Raw circle & 63.764 & 64.729 & 0.986 & 1.003 & 1.011 \\
Optimized positive & 92.861 & 91.446 & 1.015 & 1.002 & 1.003 \\
Early tree positive & 37.904 & 38.139 & 0.992 & 0.951 & 1.020 \\
Late tree positive & 140.885 & 136.867 & 1.028 & 1.065 & 1.007 \\
New face positive & 120.063 & 206.335 & 0.582 & 0.955 & 0.973 \\
Trefoil & 41.078 & 64.424 & 0.638 & 0.984 & 1.025 \\
Genus-one miss & 43.752 & 66.204 & 0.667 & 0.993 & 1.000 \\
\bottomrule
\end{tabular}

\end{center}

The new seven-crossing native positive takes about 206 milliseconds with
face search versus 120 milliseconds when the existing exact fallback
finishes the call. Its paired ratio corresponds to about 72 percent
additional time. The two continuing misses cost about 50 and 57 percent
more. Earlier positive branches avoid the new search; their small timing
differences provide no evidence of algorithmic gains and are accompanied
by visible A/A variation. In particular, the late-tree off-arm A/A ratio
is 1.065. The experiment therefore supports optional discovery and stronger
coverage evidence, not a faster default recognizer. The maintained default
continues to request zero face roots.

Reproduction uses \path{../fast/normal_orbit_research/face.py} with
\code{audit} or \code{recognize}; the audit requires optional Regina.
The existing incoming archive inventory is unchanged. The checked
\code{openai/math} head remains \code{fd4aeeb2e}, and the arXiv record for
Lackenby's hierarchy paper~\cite{lackenby2026} still lists its July 2026 first version. The
new optimal-face argument is a local extension of report 31's matching
duality, not a resolution of the general recognition-complexity gap.
